\section{A fourth-moment check near one}
This independent bound gives a short check of the analytic threshold
theorem. It is not used in the main paper's proofs.

\begin{theorem}[two modes for memory near one]\label{thm:two-modes-near-one}
If $\frac45\le a<1$, that is $0<\varepsilon\le\frac1{10}$, then case (b)
of Theorem~\ref{thm:two-modes} holds. So $g_a$ has exactly two modes.
\end{theorem}

\begin{proof}
In case (a), Khinchine's representation~\cite{dj1988} gives
$\mathcal L^\circ_a\overset{d}=UZ$, with $U$ uniform on $(0,1)$ and
independent of $Z$. Hence Cauchy--Schwarz gives
$\E(\mathcal L^\circ_a)^4\ge\frac95[\E(\mathcal L^\circ_a)^2]^2$.
The moment recursion of~\cite[Theorem~1.4]{glr2026} gives the actual ratio
\[
 \frac{\E\mathcal L_a^4}{(\E\mathcal L_a^2)^2}
 =\frac{6(2a^2+2a-1)}{4a-1}\,\operatorname{Beta}(2a,2a),
\]
where $\operatorname{Beta}$ is the beta function. For $a\ge\frac34$, its integral and
$t(1-t)\le\frac14$ give
$\operatorname{Beta}(2a,2a)\le4^{-2a+3/2}\operatorname{Beta}(\frac32,\frac32)=\pi2^{-4a}$.
Put $R(a)=(2a^2+2a-1)/(4a-1)$. On $[\frac34,1]$ we have
$(\log R)'\le2<4\log2$. The first inequality reduces to
$4a(4a^2+a-2)\ge0$. Thus $6\pi R(a)2^{-4a}$ decreases there, and for
$a\ge\frac45$ the displayed moment ratio is at most
\[
 \frac{282\pi}{55\,2^{16/5}}
 <\frac{141}{80}<\frac95,
\]
using $\pi<22/7$ and $2^{1/5}>8/7$. Even moments are unchanged by
symmetrization, so this contradicts case (a). Theorem~\ref{thm:two-modes}
then gives case (b).
\end{proof}
